\begin{table}[htbp]
\centering
\small
\caption{The Industrial Crowding-Out Effect: Impact of Power Tariff Shocks on Manufacturing}
\label{tab:table3_crowdout}
\adjustbox{max width=\textwidth}{%
\begin{tabular}{lcccc}
\toprule
 & \multicolumn{3}{c}{$\ln(\text{Employment})$: NAICS 325 \& 331} & $\ln(\text{Mfg Weekly Wage})$ \\
\cmidrule(lr){2-4} \cmidrule(lr){5-5}
 & (1) & (2) & (3) & (4) \\
\midrule
Industrial Tariff (cents/kWh) & -0.0412*** & -0.0408*** & -0.0325*** & -0.0142*** \\
 & (0.0084) & (0.0084) & (0.0088) & (0.0032) \\
Industrial Tariff $\times$ Constrained LDA & -- & -- & -0.0189*** & -- \\
 & & & (0.0051) & \\
$\ln(\text{Data Center MW} + 1)$ & -- & -0.0085** & -0.0072* & -0.0028** \\
 & & (0.0038) & (0.0039) & (0.0012) \\
\midrule
County Fixed Effects & Yes & Yes & Yes & Yes \\
Quarter Fixed Effects & Yes & Yes & Yes & Yes \\
Observations & 3,600 & 3,600 & 3,600 & 3,600 \\
$R^2$ & 0.935 & 0.936 & 0.941 & 0.928 \\
\bottomrule
\end{tabular}%
}
\begin{minipage}{\linewidth}
\vspace{1ex}
\footnotesize \textit{Notes:} Standard errors clustered at the county level in parentheses. * $p<0.10$, ** $p<0.05$, *** $p<0.01$. Data covers energy-intensive, trade-exposed manufacturing (NAICS 325 Chemical Manufacturing and NAICS 331 Primary Metal Manufacturing) from the BLS QCEW. Column (1) documents a statistically significant semi-elasticity of -0.0412 ($t = -4.90$), indicating a 4.12\% employment contraction per 1 cent/kWh tariff increase. Column (3) reveals a negative interaction wedge (-0.0189, $t = -3.71$), confirming accelerated crowding out in transmission-constrained zones.
\end{minipage}
\end{table}
